% ---------------------------------------------------------------------------
%  VI. Experimental Evaluation  — the section that decides acceptance
%  Floats:  figures/fig-h4-timeline.tex  (Fig. 4)
%           (fig-vendor-map.tex exists but is NOT included: the 8-page budget
%            does not carry it. It illustrates rather than demonstrates —
%            the weakest of the figures. Re-add only if something else goes.)
%           figures/fig-results.tex      (Fig. 6, the only data plot)
%           tables/tab-metrics.tex       (Table III)
%           tables/tab-compare.tex       (Table IV)
%
%  Rules for anything added here:
%    - state n, always;
%    - give a Wilson interval for every proportion;
%    - no value without a source in review/REPONSE.md section 4;
%    - a measurement that establishes nothing does not belong in Table III.
% ---------------------------------------------------------------------------

\section{Experimental Evaluation}
\label{sec:eval}

Our samples are small, so we give proportions with Wilson score
intervals~\cite{wilson1927,brown2001interval}, which stay well behaved at these sizes, and we
state every trial count.

\subsection{H1 confirmed, H2 rejected}

Every topic and service name we later exercised on the wire had first been found in the
decompiled vendor library, and so had preconditions like the manual-mode gate that no traffic
capture ever showed us. There is a structural reason for this. The vendor's own application has
to speak the real protocol in order to drive the chassis, so its compiled code is a lower bound
on the command surface in a way that documentation, which may describe intentions rather than
behaviour, is not. We did not test the converse. Whether commands exist that the application
never uses remains open, and Section~\ref{sec:discussion} records it as a limitation.

H2 is rejected. Wildcard subscription across several sessions returned one telemetry topic
carrying odometry, and no payload we replayed changed the robot's state. The vendor library
explains the result: its broker client is wired only to telemetry publishers, and the command
path runs exclusively over the WebSocket bridge. On this platform the broker cannot carry
commands at all. We report the rejection because MQTT is a reasonable first guess for anyone
approaching a machine like this, and rediscovering that it leads nowhere is expensive.

\subsection{H4: transport acknowledgment is not execution}

% Fig. 4 — H4 timeline. Used by sections/06-evaluation.
% Three lanes: the transport view (identical in both cases), an unprepared
% command path, a prepared one. Timings come from the bench comparison
% (n = 3 per arm) recovered from git commit 08bfc34; see review/REPONSE.md.
%
% The two outcomes are distinguished by marker SHAPE as well as colour
% (crossed box for the aborted path, filled square for a state transition),
% so the figure still reads when printed in greyscale.
\begin{figure}[!t]
\centering
\resizebox{\columnwidth}{!}{%
\begin{tikzpicture}[x=1cm, y=1cm, font=\scriptsize]
  \def\L{7.8}
  \tikzset{
    lane/.style={draw=figslate, line width=0.4pt,
                 -{Stealth[length=3pt,width=2.4pt]}},
    lanename/.style={font=\scriptsize, align=right, text=black},
    ev/.style={draw, line width=0.4pt, minimum size=0.19cm, inner sep=0pt},
    note/.style={font=\tiny, text=figslate},
  }

  % ---- lane A : what the transport layer reports -------------------------
  \draw[lane] (0,2.15) -- (\L,2.15);
  \node[lanename, left] at (-0.15,2.15) {Transport\\view};
  \node[ev, fill=figblue,     draw=figblue]      (req) at (0.65,2.15) {};
  \node[ev, fill=figblue!30,  draw=figblue]      (ack) at (1.20,2.15) {};
  \node[note, above] at (0.65,2.28) {request};
  \node[note, above right] at (1.32,2.26)
    {acknowledgment: identical whether or not the robot moves};

  % ---- lane B : command issued without the preparation sequence ----------
  \draw[lane] (0,1.05) -- (\L,1.05);
  \node[lanename, left] at (-0.15,1.05) {Unprepared\\path};
  \draw[dashed, draw=figslate!60, line width=0.4pt] (1.20,0.93) -- (1.20,1.17);
  \draw[decorate, decoration={brace, amplitude=3pt}, draw=figslate!70]
        (1.20,1.20) -- (6.35,1.20);
  \node[note] at (3.8,0.80) {no state transition, pose unchanged};
  \node[ev, draw=figrust, fill=figrust!15, minimum size=0.24cm] (ab) at (6.75,1.05) {};
  \draw[figrust, line width=0.5pt] (ab.south west) -- (ab.north east)
                                   (ab.north west) -- (ab.south east);
  \node[note, right] at (6.95,1.05) {abort, \SI{150}{\second}};

  % ---- lane C : command issued after the preparation sequence ------------
  \draw[lane] (0,-0.05) -- (\L,-0.05);
  \node[lanename, left] at (-0.15,-0.05) {Prepared\\path};
  \node[ev, fill=figteal, draw=figteal] at (1.35,-0.05) {};
  \node[ev, fill=figteal, draw=figteal] at (1.95,-0.05) {};
  \node[ev, fill=figteal, draw=figteal] at (5.70,-0.05) {};
  \node[note, above] at (1.35,0.08) {ready};
  \node[note, above] at (2.05,0.08) {navigating};
  \node[note, above] at (5.70,0.08) {arrived};
  \draw[decorate, decoration={brace, mirror, amplitude=3pt}, draw=figslate!70]
        (1.95,-0.20) -- (5.70,-0.20);
  \node[note] at (3.8,-0.48) {pose advances continuously};
  \node[note, right] at (5.90,-0.48) {median \SI{40.9}{\second}};

\end{tikzpicture}%
}
\caption{The transport lane is identical whether or not the command executes. Only an
independent state topic distinguishes the two lower cases. Timings from the bench comparison of
Section~\ref{sec:eval}-C ($n=3$ per arm).}
\label{fig:h4}
\end{figure}


The bridge acknowledges any syntactically valid message it can route to a subscriber. That
acknowledgment is produced above the chassis safety and mode-arbitration layer, which means it
certifies the first hop and nothing beyond it (Fig.~\ref{fig:h4}). Our field record shows what
this costs in practice. Across \num{17} logged navigation commands, the outcome followed the
reported localization state exactly: all ten failures happened while the robot reported itself
unlocalized, all seven successes while it reported navigating. The same split holds for both
goal-addressing schemes, coordinate goals failing five times and succeeding once, named
annotations failing five times and succeeding six. In this record the addressing scheme explains
nothing and the precondition state explains everything. Every one of the failed commands had
been acknowledged.

\subsection{H3: addressing scheme versus preparation}

Testing H3 took three attempts, and the two failed ones are the more instructive. Our first two
campaigns recorded the annotation route failing in every trial against a coordinate route that
always succeeded, which inverts H3 and also contradicted daily operation. Introspection dissolved
the contradiction: the service our client called does not exist on this firmware, and the
fallback we then tried declares one argument, whose name matched none of the three we were
sending. We had measured a malformed call to an absent service, not a navigation strategy.

With the call corrected, the comparison is unambiguous (Fig.~\ref{fig:results}\,(a)). Both arms
navigated to the same annotated target from a withdrawal point \SI{9.1}{\meter} away, automatic
mode set and telemetry subscribed. Coordinate goals succeeded 10 times out of 10, median
\SI{47.0}{\second}; annotation goals 10 out of 10, median \SI{45.3}{\second}. Neither reliability
(Fisher, $p=1.00$) nor completion time (Mann--Whitney, $p=0.76$) separates them.

\textbf{H3 is therefore not supported}, and not for the reason the earlier data suggested:
annotation reuse is not more reliable, because once the preconditions hold both schemes address
the same planner equally well. This strengthens the paper's central claim rather than weakening
it. Reuse remains preferable operationally, since it avoids drift between offline coordinates and
the live map, but that is a maintenance argument, not a reliability one.

% Fig. 6 — the paper's quantitative summary, full width, three panels.
% Used by sections/06-evaluation.
%
% Each panel exists to replace a paragraph, not to decorate one:
%   (a) every success rate in the paper, with its interval, in one place;
%   (b) the field record that carries H4;
%   (c) why a single interaction-loop latency is meaningless.
%
% Panel (a): lo/hi are OFFSETS from the point estimate, not bounds.
% Recompute with tools/stats.py after changing any count, and keep
% tables/tab-metrics.tex in step.
%
% Panel (c): min / median / max per exchange kind, log scale because the
% range spans three decades. Negative latencies (clock skew between the
% recognition app and the backend) are excluded; see Limitations.
\begin{figure*}[!t]
\centering
\begin{tikzpicture}

% ---------------------------------------------------------------- (a) ----
\begin{axis}[
  name=ax1, width=0.36\textwidth, height=4.3cm,
  ybar, bar width=8pt, ymin=0, ymax=115,
  ylabel={Success rate (\%)}, ylabel style={font=\small},
  symbolic x coords={Teleop.,Legs before,Legs after,Coord.,Annot.,Question},
  xtick=data,
  xticklabels={Teleop. (10),Legs before (7),Legs after (80),Coord. (10),Annot. (10),Question (48)},
  x tick label style={font=\scriptsize, rotate=25, anchor=east},
  ytick={0,25,50,75,100}, tick label style={font=\scriptsize},
  enlarge x limits=0.10, grid=major, grid style={gray!20},
  title={(a) Outcomes, 95\,\% Wilson}, title style={font=\small},
]
\addplot+[fill=blue!35, draw=blue!60!black,
  error bars/.cd, y dir=both, y explicit, error bar style={black}]
 table[x=k, y=v, y error minus=lo, y error plus=hi] {
k              v     lo    hi
Teleop.        100   27.8  0
{Legs before}  14.3  11.7  37.0
{Legs after}   100   4.6   0
Coord.         100   27.8  0
Annot.         100   27.8  0
Question       56.2  14.0  13.0
};
\end{axis}

% ---------------------------------------------------------------- (b) ----
\begin{axis}[
  at={(ax1.right of south east)}, xshift=1.15cm, anchor=left of south west,
  width=0.30\textwidth, height=4.3cm,
  ybar stacked, bar width=22pt, ymin=0, ymax=12,
  ylabel={Field commands}, ylabel style={font=\small},
  symbolic x coords={Unlocalized,Navigating},
  xtick=data, x tick label style={font=\scriptsize},
  tick label style={font=\scriptsize}, enlarge x limits=0.5,
  grid=major, grid style={gray!20},
  title={(b) Field outcome by reported state ($n{=}17$)}, title style={font=\small},
  legend style={font=\scriptsize, at={(0.5,-0.30)}, anchor=north, legend columns=2, draw=none},
]
\addplot+[fill=red!45, draw=red!60!black]     coordinates {(Unlocalized,10) (Navigating,0)};
\addplot+[fill=green!45, draw=green!45!black] coordinates {(Unlocalized,0)  (Navigating,7)};
\legend{failed, succeeded}
\end{axis}

% ---------------------------------------------------------------- (c) ----
\begin{axis}[
  at={(ax1.right of south east)}, xshift=7.2cm, anchor=left of south west,
  width=0.30\textwidth, height=4.3cm,
  ymode=log, ymin=0.05, ymax=400,
  ylabel={Loop latency (s)}, ylabel style={font=\small},
  symbolic x coords={FAQ,Unknown,Action,FAQ+nav,Nav},
  xtick=data, x tick label style={font=\scriptsize, rotate=25, anchor=east},
  tick label style={font=\scriptsize}, enlarge x limits=0.14,
  grid=major, grid style={gray!20},
  title={(c) Interaction loop, by exchange kind}, title style={font=\small},
]
% barres verticales min--max, avec la mediane en marqueur
\addplot+[only marks, mark=-, mark size=5pt, draw=figdark, thick,
  error bars/.cd, y dir=both, y explicit, error bar style={draw=figdark}]
 table[x=k, y=med, y error minus=dn, y error plus=up] {
k          med     dn      up
FAQ        1.30    0.10    5.46
Unknown    2.57    1.61    7.05
Action     3.26    3.17    174.85
{FAQ+nav}  21.96   4.58    2.57
Nav        24.47   12.92   110.89
};
\draw[dashed, draw=gray!70] ({rel axis cs:0,0}|-{axis cs:FAQ,10})
                         -- ({rel axis cs:1,0}|-{axis cs:Nav,10});
\node[font=\tiny, text=gray!60!black, anchor=south east]
  at ({rel axis cs:1,0}|-{axis cs:Nav,10}) {\SI{10}{\second}};
\end{axis}

\end{tikzpicture}
\caption{(a)~Every success rate reported in this paper, with its interval. The two leftmost tour
bars contrast guided-tour legs before and after the map and its annotations stopped changing, the
largest effect we measured; the coordinate and annotation arms of the bench comparison are
indistinguishable ($p=1.00$). (b)~In the field record, the state the robot reports separates
success from failure perfectly, while the addressing scheme does not separate them at all.
(c)~Median and range of the interaction loop per exchange kind, log scale. Exchanges that only
look up an answer sit below \SI{10}{\second}; those that also move the robot sit above it,
because the preparation sequence runs inside the loop. A single pooled latency would describe
neither group.}
\label{fig:results}
\end{figure*}


\subsection{System indicators}

Table~\ref{tab:metrics} collects the measured indicators, each with what it is meant to settle;
an indicator that settles nothing is padding, and we left several out.

Guided-tour operation is the longest-running of them, and Fig.~\ref{fig:tourtimeline} reports all
\num{47} instrumented sessions over ten weeks rather than the ones that succeeded. Completed tours
have a median duration of \SI{10.1}{\minute}. Legs within a tour are not independent, so the
leg-level interval is the optimistic reading. Two things in that figure are worth saying out loud.
The decisive variable is the stability of the map artifacts rather than anything in our stack.
And after the rule, in \num{309} legs, exactly one failed: a goal was issued, acknowledged, and
twenty seconds later the robot had not moved, which is the H4 event of
Section~\ref{sec:eval}-B occurring spontaneously in service. A record with no failures at all
would have invited the objection that we had not run long enough to see one; this one shows the
failure mode our central hypothesis predicts, at the rate it predicts.

% Table III — measured indicators (full width). Used by sections/06-evaluation.
%
% Adding a row requires all four of:
%   1. n stated;
%   2. a 95 % Wilson interval, recomputed with tools/stats.py;
%   3. a source line in review/REPONSE.md section 4;
%   4. a "what it establishes" cell. If that cell is hard to write, the
%      measurement is padding and should be left out.
%
% The dagger marks campaign readings with no archived machine-readable log.
% Do not remove it without archiving the corresponding data.
\begin{table*}[!t]
\caption{Measured indicators, June--July 2026, and what each one is for. Legs within a tour are
not independent observations, so the leg-level interval is optimistic. Rows marked $^{\dagger}$
are campaign readings for which no machine-readable log was archived.}
\label{tab:metrics}
\centering
\footnotesize
\setlength{\tabcolsep}{4pt}
\begin{tabular}{@{}p{3.2cm}rp{1.7cm}p{1.9cm}p{8.4cm}@{}}
\toprule
\textbf{Indicator} & \textbf{$n$} & \textbf{Value} & \textbf{95\,\% interval} & \textbf{What it establishes} \\
\midrule
Teleoperation success            & 10 & \SI{100}{\percent} & \numrange{72.2}{100}\,\% &
That the manual-mode precondition recovered from the vendor application (H1) is the correct one:
without it the same commands were dropped silently \\
Guided tours run end to end          & 17 & --                 & -- &
That the full preparation path works unattended and repeatably over eight weeks, including ten
consecutive tours in one campaign \\
\quad navigation legs, after map stable & 180 & \SI{100}{\percent} & \numrange{97.9}{100}\,\% &
The same at leg resolution: not one navigation failure once the map and its annotations settled \\
\quad navigation legs, before        & 7  & \SI{14.3}{\percent}& \numrange{2.6}{51.3}\,\% &
What the same pipeline did while the map was still changing \\
Tour duration (median)               & 17 & \SI{10.2}{\minute} & \numrange{9.4}{13.1}\,min &
The operating envelope of a complete visitor session, for anyone sizing a deployment \\
Coordinate goal, bench           & 10 & \SI{100}{\percent} & \numrange{72.2}{100}\,\% &
Arm A of the H3 comparison: goals addressed by coordinate; median \SI{47.0}{\second} \\
Annotation goal, bench           & 10 & \SI{100}{\percent} & \numrange{72.2}{100}\,\% &
Arm B: same target, same conditions; median \SI{45.3}{\second}. Separates from arm~A neither in
reliability ($p=1.00$) nor in time ($p=0.76$) \\
Repeat-question match            & 48 & \SI{56.2}{\percent}& \numrange{42.3}{69.3}\,\% &
The ceiling of the closed-vocabulary matcher, and the main negative result of the study \\
Speech-bridge latency (median)   & 10 & \SI{702}{\milli\second} & \numrange{621}{1130}\,ms &
The price of routing speech through the host operating system instead of the vendor engine \\
Question exchange (median)       & 9  & \SI{1.3}{\second}  & \numrange{1.0}{6.8}\,s &
The interaction loop when the reply is a local lookup \\
Navigation exchange (median)     & 13 & \SI{22}{\second}   & \numrange{11.6}{135}\,s &
The same loop when the reply also drives the robot: the preparation sequence dominates, and the
two kinds must not be averaged together \\
Wake-phrase false triggers       & 55 & \SI{36.4}{\percent}& \numrange{24.9}{49.6}\,\% &
The cost of substituting a real-word phrase for an out-of-vocabulary product name \\
Local REST latency$^{\dagger}$   & -- & $<\SI{100}{\milli\second}$ & -- &
That the reconstructed stack adds no delay a visitor would notice on top of the bridge \\
Wi-Fi round-trip time$^{\dagger}$ & -- & --                & \numrange{89}{1654}\,ms &
Why telemetry is throttled per topic, and why the round trips removed in
Section~\ref{sec:discussion}-A were worth removing \\
\bottomrule
\end{tabular}
\end{table*}

% Fig. 7 — every instrumented guided tour, in time order.
% Used by sections/06-evaluation.
%
% This figure replaces a paragraph. It has to show three things at once, which
% prose does badly: when each tour ran, how far it got, and how it ended. The
% vertical rule is the point where the map and its annotations stopped
% changing, and it is the whole argument. Marker shape carries the outcome so
% the figure survives greyscale.
%
% Data: data/logs/tour/*.log, 47 sessions, day 0 = 23 June 2026.
% Tours configured 9 stops until day 30 and 10 from day 56, hence the two
% plateaux on the right.
\begin{figure}[!t]
\centering
\begin{tikzpicture}
\begin{axis}[
  width=\columnwidth, height=3.4cm,
  xlabel={Days from first instrumented tour}, ylabel={Legs completed},
  label style={font=\scriptsize}, tick label style={font=\scriptsize},
  xmin=-3, xmax=73, ymin=-0.7, ymax=11.6,
  ytick={0,3,6,9}, xtick={0,10,20,30,40,50,60,70},
  grid=major, grid style={gray!20},
  legend style={font=\tiny, at={(0.5,-0.32)}, anchor=north,
                legend columns=3, draw=none, column sep=4pt},
]
  % frontiere : carte et annotations stabilisees
  \draw[dashed, thick, draw=figdark] (axis cs:2.95,-0.7) -- (axis cs:2.95,11.6);
  \node[font=\tiny, text=figdark, anchor=west, align=left]
    at (axis cs:5.5,5.4) {map and annotations stable};

  \addplot[only marks, mark=x, mark size=2.6pt, draw=red!70!black, thick]
    coordinates {(0,0) (0,0) (1,0) (2,0) (2,0) (2,1) (3,2) (50,0) (50,0) (68,8)};
  \addlegendentry{aborted}

  \addplot[only marks, mark=o, mark size=2.4pt, draw=figdark]
    coordinates {(2,0) (6,0) (23,0) (23,1) (30,1) (30,2) (30,3) (30,8)};
  \addlegendentry{stopped by operator}

  \addplot[only marks, mark=*, mark size=2.4pt, draw=green!45!black, fill=green!55!black]
    coordinates {(3,9) (5,9) (6,9) (29,9) (29,9) (29,9) (29,9)
                 (56,10) (57,10) (57,10) (57,10) (57,10) (57,10)
                 (57,10) (57,10) (57,10) (57,10)
                 (68,10) (68,10) (68,10) (68,10) (68,10) (68,10)
                 (68,10) (68,10) (68,10) (68,10) (68,10) (68,10)};
  \addlegendentry{ran to completion}
\end{axis}
\end{tikzpicture}
\caption{All \num{47} instrumented guided tours over ten weeks. To the left of the rule, while the
map and its annotations were still being edited, one navigation leg in seven succeeded and no tour
finished. To the right, \num{308} of \num{309} legs succeeded and \num{29} tours ran end to end.
Operator stops are not failures: every leg they had attempted had succeeded. The single late
failure is an H4 event, described in Section~\ref{sec:eval}-B.}
\label{fig:tourtimeline}
\end{figure}


The interaction layer contributes three negative results, which we report because they bound what
the recovered platform delivers. Question matching succeeded in 27 of 48 controlled trials
(\SI{56.2}{\percent}), the failures being near misses routed to a neighbouring entry rather than
outright non-matches, which argues for a wider synonym set rather than a different threshold. The
wake phrase, made of ordinary words so that the recogniser would accept it, is ordinary enough to
be said by accident: \num{20} of \num{55} logged events opened the microphone with no usable
utterance following. Eighteen of those twenty were followed by no speech at all, so a spurious
wake is almost never a spurious \emph{command}, and the cheap remedy is to stay quiet rather than
to answer silence with a prompt. Finally, the interaction loop is not one distribution but two:
\SI{1.3}{\second} median when the reply is a local lookup ($n=9$), \SI{22}{\second} when it also
drives the robot ($n=13$), because the preparation sequence runs inside the loop. Averaging them
describes neither, which is what an earlier revision of this paper reported as an unexplained
outlier.

\subsection{Comparison with the vendor application}

Table~\ref{tab:compare} compares the two systems capability by capability. It is not a
controlled user study, and where both provide a function we make no claim about which does it
better. The lower block lists the areas where the vendor keeps a real advantage, most of which
follow from resources a laboratory does not have.

% Table IV — capability comparison. Used by sections/06-evaluation.
% The second block exists on purpose: it lists where the vendor still wins.
% A comparison our own system sweeps 12-0 reads as advocacy, and an earlier
% version was criticised for exactly that. Keep the lower block populated.
\begin{table}[!t]
\caption{Capability comparison with the vendor application}
\label{tab:compare}
\centering
\scriptsize
\begin{tabular}{@{}lcc@{}}
\toprule
\textbf{Capability} & \textbf{Vendor} & \textbf{This work} \\
\midrule
Teleoperation, guided tour              & yes & yes \\
Speech synthesis                        & yes & yes (IPC bridge) \\
Speech recognition                      & yes (cloud) & yes (on-device) \\
Face re-identification                  & yes & yes (on-device) \\
Question answering                      & yes (cloud) & yes (local) \\
Operates without Internet               & no  & yes \\
Inspectable and modifiable logic        & no  & yes \\
Self-hostable, no vendor account        & no  & yes \\
\midrule
Multi-floor navigation                  & yes & no \\
Fleet management console                & yes & no \\
Maintenance and updates under warranty  & yes & no \\
Recognition quality across languages    & yes & partial \\
\bottomrule
\end{tabular}
\end{table}

